\section{Scope \& Setup}
\label{sec:scope}

\paragraph{Layer-level scope.}
This paper separates two layers that the surrounding program treats as independent.
The \emph{Lean Proof layer} establishes semantic statements at the axiom level using polymorphic Lean definitions of \(\mathtt{UN}\), \(\mathtt{MINLIB}\), \(\mathtt{Decisive}\), and \(\mathtt{SocialAcyclic}\).
The \emph{CNF Witness/Audit layer} establishes finite-case satisfiability or unsatisfiability via clause-level encodings whose conformance to the Proof-layer axioms must itself be audited.
The main positive deliverable belongs to the Proof layer; the main negative deliverable concerns the wall between the two layers.

\paragraph{Family scope.}
The positive theorem ranges over all finite cases \((\mathtt{Fin}\ n, \mathtt{Fin}\ m)\) with \(2 \leq n\) and \(4 \leq m\).
It does not claim a result for infinite-domain social choice; it does not claim a result for Arrow's theorem; and it does not claim a result that bypasses \(\mathtt{UN}\) or \(\mathtt{MINLIB}\).
The phrase ``Sen's impossibility lifts'' in this paper always means the semantic, finite-case lift in the sense above.

\paragraph{Lever scope.}
The CNF Witness/Audit layer in this repository uses an explicit lever schema in which rationality is approximated by the local short-cycle levers \(\mathtt{no\_cycle3}\) and \(\mathtt{no\_cycle4}\).
At the base case \((n,m)=(2,4)\), \(\mathtt{no\_cycle3} \wedge \mathtt{no\_cycle4}\) is equivalent to full acyclicity on the social strict preference, and that equivalence is audited exhaustively by \code{scripts/check\_acyclicity\_short\_cycles.py}.
At larger cases the audited equivalence does not transfer automatically; the existing interpretation notes record that \(\mathtt{no\_cycle3} \wedge \mathtt{no\_cycle4}\) still permits length-5 cycles at \(m=5\).
The negative-side scope wall is exactly the statement that the Proof-layer lift does not by itself extend this audited equivalence.

\paragraph{What is in this paper, what is not.}
The paper contains
(i) the polymorphic Lean signatures used by the lift theorem,
(ii) the kernel-checked statement and proof sketch of \(\mathtt{sen\_impossibility\_lifts}\),
(iii) the Proof-vs-Witness/Audit scope wall, and
(iv) the artifact binding for both deliverables.
The paper does \emph{not} contain
(i) new solver experiments,
(ii) any modification of the sibling Evidence-Contract or Explanation-Contract manuscripts,
(iii) a family-level CNF certificate for full acyclicity, or
(iv) a family-scale representation-canonicity result.
These last two are explicitly deferred to a companion paper (in preparation).
